Table~\ref{tab:ablation} removes one component of the protocol at a time and reports, on the same data, the quantity that the component makes visible, with the component and without it. The ablation is one at a time: it does not measure interactions between components.

\begin{table}[htbp]
\centering
\setlength{\tabcolsep}{3.5pt}
\footnotesize
\caption{Ablation of the proposed protocol: each row removes one component and reports, on the same data, the quantity the component makes visible; the first two rows read the committed CEREP artifacts, the next four the known-truth draws, the last two the exact enumeration and the matrix tests.}
\label{tab:ablation}
\resizebox{\linewidth}{!}{\begin{tabular}{p{3.2cm}p{2.6cm}p{3.2cm}p{2.4cm}p{3cm}}
\toprule
Component of the protocol & Without it & Quantity & With the component & Without the component \\
\midrule
\relax Report the three pages' definitions side by side & \relax Report the transition page only & \relax S\&P CCC to C, long-run PD (\%) & $26.64$ and $8.61$ & $8.61$ (the gap is not seen) \\
\relax The same, carried to capital & \relax Report the transition page only & \relax S\&P risk weight (\%) & $58.0$ and $47.5$ & $47.5$ ($10.5$ points lower) \\
\relax Informativeness ladder of $k$ & \relax Run $k=0$ only & \relax B, removed estimator, relative bias (\%) & $-35.1\pm0.3$ at $k=9$ & $2.6\pm0.4$ at $k=0$ (opposite sign) \\
\relax Known-truth reference & \relax Compare the three treatments with one another & \relax B at $k=9$: bias of the least biased treatment (followed) & $-15.6\pm0.3$ against the truth & \relax treatments span $23.7$ points; the offset of every one from the truth is unknown \\
\relax First-year cohort alone & \relax Pool five annual cohorts & \relax B, followed estimator, relative bias at $k=9$ (\%) & $1.4\pm0.8$ (first year) & $-15.6\pm0.3$ (pooled) \\
\relax Paired observed and latent draws & \relax Independent draws & \relax Standard error of removed minus latent, B, $k=9$ ($10^{-5}$) & $3.0$ & $5.6$ \\
\relax Exact enumeration of the coverage & \relax Estimate the coverage by simulation & \relax Repetitions to resolve the Wald coverage of AAA at $n=50$ to 10\% relative error & \relax 0 (exact) & $30{,}400$ \\
\relax Matrix tests as the screen for withdrawals & \relax None needed (no screen) & \relax Rejection rate at 5\% of the weaker test at $k=9$ (\%) & $1.4$ and $5.8$ & \relax the screen rejects at about its size, and the bias is not seen \\
\bottomrule
\end{tabular}}
\end{table}


The first two rows concern the definition-explicit report. A validator who reads only the transition page sees $\spLraDfourCcc\%$ for CCC to C where the default page shows $\spLraDtwoCcc\%$, and a risk weight of $\spRwDfour\%$ where the default page gives $\spRwDtwo\%$, $\spRwGap$ points lower. Nothing in the transition page reveals the difference, because the counts that would reveal it are on another page. The component costs nothing to run, since the pages are published together, and it supplies the largest single effect of the protocol on capital.

The third row concerns the ladder of informativeness. At $k=0$ the removal treatment overstates the PD of B by $\removedBkZero\pm\removedBkZerose$ percent, a result from which a validator would conclude that the industry convention is slightly conservative; at $k=9$ the same estimator understates the PD by $\removedBkNine\pm\removedBkNinese$ percent. The sign changes between $k=0$ and $k=1$. A measurement at one level of $k$ gives an answer that depends on the level chosen, and the ladder shows that dependence.

The fourth row concerns the reference to the known truth. Without the truth, the three treatments can only be compared with one another; at $k=9$ they span $\keptBkNine\%$ to $\followedBkNine\%$ and the one with the smallest bias, the followed treatment, would be taken as the reference, although it is itself below the truth by $\followedBkNine\pm\followedBkNinese$ percent. The error of the best treatment is invisible without the truth.

The fifth row concerns pooling. The followed treatment is exact on the first cohort ($\followedYoneBkNine\pm\followedYoneBkNinese$ percent at $k=9$) and biased by $\followedBkNine\pm\followedBkNinese$ percent when five cohorts are pooled, which is the selection of Proposition~\ref{prop:later}. Pooling is the practice that the long-run average of paragraph 84 of the guidelines invites, and the ablation shows the cost for the treatment that otherwise has none.

The sixth row concerns the pairing of the observed and latent draws. The standard error of the contrast between the removal estimator and the reference at $k=9$ for grade B is $\seRemovedPaired\times10^{-5}$ with pairing and $\seRemovedUnpaired\times10^{-5}$ without, a ratio of $\seRatio$, so pairing removes $\seRatioPct\%$ of the standard error of the contrast, or equivalently the pairing is worth about $\seRatio^2$ times as many repetitions. The seventh row concerns the enumeration. To resolve a Wald coverage of $\waldAaaFifty\%$ to a relative error of 10\% by simulation takes about $\repsNeeded$ repetitions, whereas the enumeration is exact; the more important gain is that a simulation of $\repsNew$ repetitions cannot distinguish a coverage of $\waldAaaFifty\%$ from zero. The eighth row concerns the screen: the matrix tests reject $\testthkNinefive\%$ and $\testmrkNinefive\%$ of the repetitions at $k=9$, no more than their size, so they would add nothing to a decision about withdrawals there, and the known-truth measurement is what reveals the bias.
